\section{Exhaustion of the nonbinary orbit actions}
\label{sec:nonbinary-alphabet}

\begin{theorem}[Complete nonbinary action reduction]
\label{thm:nonbinary-alphabet}
Let $\mathcal X_n$ consist of the actual subgroups of $S_n$ with
an orbit action outside the following list:
\[
 \text{binary transitive actions},\qquad
 \text{natural }C_3,A_4,S_3,\qquad\text{fixed points}.
\]
There are positive constants $c,c'$ and complete kernels $K(n,b)$,
with $K(n,b)=0$ for $b\geq n$, such that eventually
\[
 \frac{|\mathcal X_n|}{L_n}
 \leq\eta(n)+\sum_{b<n}K(n,b)\frac{s_b}{L_b},\qquad
 \sum_{b<n}K(n,b)=O(2^{-cn}),\qquad
 \eta(n)=O(2^{-c'n^2}).
\]
Every forward term retains the complete complementary subgroup and
the original permutation-action weight. Thus every subgroup outside
this counted union has only the displayed orbit actions.
\end{theorem}
\begin{proof}
We first prove that every transitive action outside the list meets
one of the proved consumers on its entire original normal menu.
This is an action-by-action exhaustion, preceding the counting union.

For a primitive action, a nonaffine socle is covered by
Theorem~\ref{thm:semisimple-compression}. An affine action of degree
at least $1024$ is covered by Theorem~\ref{thm:small-order-consumer}.
At smaller degrees at least five, the soluble and nonsoluble cases
are Theorems~\ref{thm:soluble-primitive-consumer} and
\ref{thm:nonsoluble-primitive-consumer}, with the exceptional
degree-five Frobenius action covered by
Proposition~\ref{prop:s4-f20-consumers}. At degree four the only
primitive groups are the natural $A_4,S_4$, the latter again
covered by that proposition. Degrees three and two leave precisely
the allowed natural actions and the binary group.

Now let $U$ be imprimitive. Choose a minimal nontrivial block of
size $r$, with $s$ blocks and faithful top $T\leq S_s$.
The full block stabilizer induces an exact primitive local group
$L\leq S_r$: any smaller block for that induced action and all
its translates would give a smaller block for $U$. Transport
coordinates along representatives of the block action to obtain
the actual embedding $U\leq L\wr T$. This does not impose a
splitting or full block kernel.

If $L$ is nonaffine, use Theorem~\ref{thm:semisimple-compression}.
If it is affine with $r\notin\{2,3,4,5,8,9,16\}$, use
Theorem~\ref{thm:affine-local-capacity}. For the four degrees
$5,8,9,16$, Theorem~\ref{thm:small-affine-capacity} covers every
count except its explicit finite lists; these lists are consumed,
respectively, by Theorems~\ref{thm:five-affine-exceptions},
\ref{thm:binary-affine-exceptions}, \ref{thm:nine-affine-exceptions}
and \ref{thm:binary-affine-exceptions}. Each is an entire-normal
statement on arbitrary complete sources.

For $r=3$, the exact local group is $C_3$ or $S_3$.
Theorem~\ref{thm:prime-s3-capacity} and
Proposition~\ref{prop:bounded-c3} cover all $C_3$ counts.
The same capacity theorem and Theorem~\ref{thm:bounded-s3}
cover all $S_3$ counts.

For $r=4$, the local group is natural $A_4$ or $S_4$.
Theorem~\ref{thm:four-local-capacity} leaves only its displayed
finite count sets, all at most $3072$. In the actual affine frame
put $E=U\cap(\F_2^2)^s$, $R=U/E$ and $D=R\cap C_3^s$.
For $D\ne1$, Theorem~\ref{thm:four-nonzero-core} applies.
For $D=1$, Theorem~\ref{thm:zero-ternary-capacity} covers all
counts except $4,8,16$, with counts two and three impossible.
At four blocks use Theorem~\ref{thm:degree-sixteen-nonbinary};
at eight and sixteen use Theorem~\ref{thm:zero-ternary-small-frames}.
This disposes of every exact four-point component.

Only $r=2$, $L=C_2$ remains. Its pair kernel is binary, so its
top is nonbinary whenever $U$ is nonbinary. The prime-capacity
theorem leaves exactly
\[
 \{2^a:1\leq a\leq8\}\cup
 \{3\cdot2^a:0\leq a\leq10\}.
\]
At count two the whole action would be binary. Counts three,
four and six are Proposition~\ref{prop:three-pair-actions} and
Propositions~\ref{prop:four-pairs}, \ref{prop:six-pairs}.
Count eight is Theorem~\ref{thm:degree-sixteen-nonbinary}, and
counts twelve and sixteen are Theorems~\ref{thm:twelve-pairs}
and \ref{thm:sixteen-pairs}. The remaining necessary counts are
\[
 24,32,48,64,96,128,192,256,384,768,1536,3072.
\]
Every one is even and at least 24. For each original $N\lhd U$,
Corollary~\ref{cor:pair-fusion} applies to the actual section
$K/(K\cap N)$, where $K$ is the pair kernel, with quotient
$U/(KN)$. There is no assumption that $N\leq K$. Since these
physical degrees form a fixed finite set, Theorem~\ref{thm:fusion}
covers the union of all their normal entries with a strict gap.
This proves the required exhaustion.

Finally choose one of the proved certificates on each original
action class, before inspecting its complementary group, and give
the finitely many consumer families a fixed order. For a subgroup
in $\mathcal X_n$, at least one orbit and its actual Goursat normal
entry belong to the first eligible family. Apply that family's
complete incidence sum. For a family with a hot/cold split, the
existence of any hot eligible occurrence defines its hot part;
all eligible occurrences in its remaining part are cold. Thus a
chosen cold pointing cannot hide another hot one.

The incidence factor for deleting the original $w$-point orbit is
$n!/(b!a(U))$, with $b=n-w$. It changes the occurrence factorial
once. None of the block systems, chief factors, comparator degrees
or character tuples acquires a physical support or normalizer.
The full-fibre theorems are evaluated first on the same actual
complete $J\leq S_b$, so the cold bounds multiply $s_b/L_b$ and
the hot bounds are complete scalars at $n$. Every finite family
has strict positive margin, and every unbounded family has its
proved uniform margin and subquadratic error. A finite sum of
their exponential rows and quadratic scalars gives the asserted
$K,\eta,c,c'$. No bound on $s_b/L_b$ was assumed.
\end{proof}
